\documentclass[a4paper,10pt,twoside,openany]{article}

\usepackage[lang=hebrew]{maths}
\usepackage{hebrewdoc}
\usepackage{stylish}
\usepackage{lipsum}
\let\bs\blacksquare

\setlength{\parindent}{0pt}

%%%%%%%%%%%%
% Styling %
%%%%%%%%%%%%

\usepackage{enumitem}

%%%%%%%%%%%%%
% Counters  %
%%%%%%%%%%%%%

\setcounter{section}{1}     
            
%BIBLIOGRAPHY
\usepackage[
backend=biber,
style=alphabetic,
]{biblatex}
\addbibresource{bibliography.bib} %Imports bibliography file

%%%%%%%%%%
% Title  %
%%%%%%%%%%
\title{
אלגברה ב' - נספח לגיליון תרגילים 1 \\
הגדרות וסימונים עבור מטריצות מייצגות
}
\date{}

\begin{document}
\maketitle

\begin{definition}[וקטור קואורדינטות לפי בסיס]
יהי
$V$
מרחב וקטורי ממימד
$n \in \mbb{N}$
מעל שדה
$\mbb{F}$,
ויהי
$\mcal{B} = \prs{v_1, \ldots, v_n}$
בסיס של
$V$.
כל
$v \in V$
ניתן לכתוב בצורה יחידה בתור
\[v = \alpha_1 v_1 + \ldots + \alpha_n v_n\]
עבור
$\alpha_i \in \mbb{F}$.
נגדיר את
\emph{וקטור הקואורדינטות של
$v$
כזה לפי הבסיס
$\mcal{B}$}
בתור
\[\text{.} \brs{v}_{\mcal{B}} = \pmat{\alpha_1 \\ \vdots \\ \alpha_n}\]
\end{definition}

\begin{definition}[מטריצה מייצגת]
יהיו
$V,W$
מרחבים וקטוריים ממימדים
$n,m \in \mbb{N}$
בהתאמה מעל אותו שדה
$\mbb{F}$,
ותהי
$T \in \Hom_{\mbb{F}}\prs{V,W}$
העתקה לינארית מ־%
$V$
ל־%
$W$.
יהיו
\begin{align*}
\mcal{B} = \prs{v_1, \ldots, v_n}
\end{align*}
ו־%
$\mcal{C}$
בסיסים של
$V$
ו־%
$W$
בהתאמה.

נגדיר את
\emph{המטריצה המייצגת}
$\brs{T}^\mcal{B}_\mcal{C}$
בתור
\begin{align*}
\text{.} \brs{T}^\mcal{B}_\mcal{C} \coloneqq \pmat{\vert & & \vert \\ \brs{T\prs{v_1}}_C & \cdots & \brs{T\prs{v_n}}_C \\ \vert & & \vert}
\end{align*}

אם
$V = W$
וכן
$\mcal{B} = \mcal{C}$,
נסמן
\begin{align*}
\text{.} \brs{T}_\mcal{B} \coloneqq \brs{T}^{\mcal{B}}_{\mcal{B}}
\end{align*}
\end{definition}

\begin{proposition} \label{proposition:matrix-of-composition}
יהיו
$U,V,W$
מרחבים וקטוריים סוף־מימדיים מעל שדה
$\mbb{F}$,
עם בסיסים
$\mcal{B}, \mcal{C}, \mcal{D}$
בהתאמה,
ותהיינה
$S \in \Hom_{\mbb{F}}\prs{U,V}$
ו־%
$T \in \Hom_{\mbb{F}}\prs{V,W}$.

אז
\begin{align*}
\text{.} \brs{T \circ S}^\mcal{B}_\mcal{D} = \brs{T}^\mcal{C}_\mcal{D} \brs{S}^\mcal{B}_\mcal{C}
\end{align*}
\end{proposition}

\begin{definition}[מטריצת מעבר בסיס]
יהי
$V$
מרחב וקטורי ממימד
$n \in \mbb{N}$,
עם בסיסים
$\mcal{B},\mcal{C}$,
ותהי
$\id_V$
העתקת הזהות על
$V$.
נגדיר
\begin{align*}
\text{.} M^\mcal{B}_\mcal{C} = \brs{\id_V}^\mcal{B}_\mcal{C}
\end{align*}
\end{definition}

\begin{proposition} \label{proposition:change-of-basis-matrices}
יהי
$V$
מרחב וקטורי סוף־מימדי מעל שדה
$\mbb{F}$,
ויהיו
$\mcal{B}, \mcal{C}, \mcal{D}$
שלושה בסיסים של
$V$.

אז:
\begin{enumerate}
\item $M^\mcal{B}_\mcal{C}$
מטריצה הפיכה, ומתקיים
$\prs{M^\mcal{B}_\mcal{C}}^{-1} = M^\mcal{C}_\mcal{B}$.
\item $M^\mcal{B}_\mcal{C} = M^\mcal{D}_\mcal{C} M^\mcal{B}_\mcal{D}$.
\end{enumerate}
\end{proposition}

\newpage

\begin{remark}
למטריצה
$M^{\mcal{B}}_{\mcal{C}}$
\emph{לא נקרא}
מטריצת מעבר מהבסיס
$\mcal{B}$
לבסיס
$\mcal{C}$
או להיפך.

מצד אחד, היא מקיימת
\[\text{,} \forall v \in V : M^{\mcal{B}}_{\mcal{C}} \brs{v}_{\mcal{B}} = \brs{v}_{\mcal{C}}\]
ולכן יש הקוראים לה מטריצת מעבר מהבסיס
$\mcal{B}$
לבסיס
$\mcal{C}$.

מצד שני, אם
$V = \mbb{F}^n$,
\[\mcal{E} = \prs{e_1, \ldots, e_n}\]
הבסיס הסטנדרטי של
$V$,
ונכתוב
\begin{align*}
\mcal{B} &= \mcal{E} \\
\text{,} \mcal{C} &= \prs{c_1, \ldots, c_n}
\end{align*}
נקבל כי
\begin{align*}
M^{\mcal{B}}_{\mcal{C}} c_i &= M^{\mcal{E}}_{\mcal{C}} \brs{c_i}_{\mcal{E}}
\\&= \brs{c_i}_{\mcal{C}}
\\&= e_i
\end{align*}
כלומר
\begin{align*}
\text{.} \prs{M^{\mcal{B}}_{\mcal{C}} c_1, \ldots, M^{\mcal{B}}_{\mcal{C}} c_n} = \mcal{B}
\end{align*}
לכן, יש הקוראים למטריצה
$M^\mcal{B}_\mcal{C}$
מטריצת מעבר מהבסיס
$\mcal{C}$
לבסיס
$\mcal{B}$.

כדי למנוע בלבול, במקום להתייחס למטריצה כזאת כמטריצת מעבר מבסיס זה או אחר, נכתוב תמיד מפורשות
$M^\mcal{B}_\mcal{C}$.
\end{remark}

\end{document}